% !TeX root = ../main.tex

\section{La recapitulación y el horizonte del milenio}

Interpretar la primera resurrección exige determinar primero cómo se relacionan las visiones de Apocalipsis 19 y 20.
Una lectura rigurosamente sucesiva obligaría a situar la escena de los tronos después de la derrota de la bestia y del falso profeta; una lectura recapitulativa permite entenderla como una nueva perspectiva sobre la era que precede al juicio definitivo.
La segunda opción no depende de una predisposición contra el milenio literal, sino de la manera en que Juan dispone sus imágenes.
El Apocalipsis vuelve reiteradamente sobre el conflicto entre el Cordero y sus adversarios, mostrando su curso y desenlace desde ángulos distintos.
Hendriksen hizo de estas visiones paralelas un eje de su exposición del libro, y Bauckham ha insistido en que la teología del Apocalipsis resulta inseparable de su composición simbólica. \autocite{hendriksen1998,bauckham1993}

La relación entre Apocalipsis 19:17--21 y 20:7--10 ofrece una prueba especialmente significativa.
En el primer pasaje, un ángel convoca a las aves para devorar la carne de los ejércitos derrotados; en el segundo, Gog y Magog son congregados desde los cuatro extremos de la tierra para la batalla final.
Ambas escenas reutilizan el oráculo de Ezequiel 38--39: la reunión de las naciones para asaltar al pueblo de Dios (Ez 38:4, 8--9, 16), su destrucción mediante la intervención divina (38:18--23; 39:1--6) y el banquete de las aves tras la derrota (39:17--20).
El vínculo no consiste en una identidad de todos los detalles: el convite de Ezequiel aparece explícitamente en Apocalipsis 19, mientras el nombre de Gog y Magog se reserva para el capítulo 20.
Precisamente la distribución de esas imágenes entre ambas visiones sugiere que Juan contempla el mismo conflicto escatológico desde perspectivas complementarias. \autocite{beale1999,koester2014}

La alternativa secuencial tendría que postular dos concentraciones mundiales de enemigos y dos derrotas decisivas, separadas por mil años, a partir de un mismo trasfondo profético.
Tal reconstrucción es posible como hipótesis, pero resulta menos económica que la recapitulación.
El propio Hill observa esta dificultad en el capítulo leído para el curso y remite a la controversia exegética sobre la reutilización de Gog y Magog. \autocite{hill2016}
Además, Apocalipsis 20:1--3 no describe una anulación absoluta de la actividad de Satanás, sino su restricción con respecto al engaño de las naciones para la ofensiva final.
El dragón puede perseguir a los santos durante el tiempo de la iglesia sin conseguir todavía su exterminio universal (Ap 12:12--17).
Por consiguiente, es razonable entender que al comenzar el capítulo 20 la visión retrocede hasta la era inaugurada por la victoria de Cristo, y no que prosigue necesariamente la cronología de 19:21.
La recapitulación no demuestra por sí sola qué es la primera resurrección, pero elimina el presupuesto de que deba ocurrir después de la parusía.

\section{Regeneración y estado intermedio: dos lecturas amileniales}

Descartada la necesidad de una sucesión lineal, permanece una cuestión más precisa: ¿qué experiencia de los santos recibe el nombre de \enquote{primera resurrección}?
La interpretación amilenial no ha sido monolítica.
Agustín, al comentar Juan 5:24--29, distinguió una vivificación espiritual que acontece ahora y una resurrección corporal que tendrá lugar al final.
Aplicó esa distinción a Apocalipsis 20 y entendió la primera resurrección como el tránsito de la muerte del pecado a la vida de la gracia.
Su exposición tenía un fundamento bíblico real: el Señor declara que ha llegado la hora en que los muertos oyen su voz y viven, mientras anuncia otra hora futura en que todos los sepultados saldrán de sus tumbas (Jn 5:25, 28--29). \autocite{augustin}

Pablo utiliza un lenguaje igualmente contundente.
En Efesios 2:1--6, los que estaban muertos en delitos son vivificados, resucitados y sentados con Cristo en lugares celestiales; Colosenses 3:1 presupone que los creyentes ya han resucitado con él.
Por tanto, no hay objeción teológica a llamar resurrección a una realidad espiritual.
Pero reconocer la legitimidad de ese uso neotestamentario no resuelve la referencia concreta de Apocalipsis 20:4--6.
Allí los protagonistas no aparecen como incrédulos que reciben por primera vez la vida de la fe.
Son personas que han sostenido el testimonio de Jesús, rechazado la adoración de la bestia y afrontado la muerte por su fidelidad.
La regeneración precede a su perseverancia; su ejecución no constituye la ocasión en que dejan de estar espiritualmente muertos.

Kline propuso, por eso, identificar la primera resurrección con la entrada de los creyentes muertos en el estado intermedio de gloria.
No negó que Efesios describa una resurrección espiritual presente; distinguió esa realidad soteriológica del acontecimiento que Juan contempla al ver las almas de quienes ya han dado su testimonio.
El martirio se convierte entonces en la transición paradójica de la muerte corporal a la participación manifiesta en el reinado celestial del Cordero.
En su formulación, lo que para el mundo parece muerte definitiva es para el cristiano acceso a una vida victoriosa. \autocite[370--374]{kline1975}
Johnson desarrolla la misma identificación: los mártires son la representación especialmente dramática de quienes mueren en Cristo, y la visión contempla su vindicación ante el trono celestial. \autocite{johnson2001,johnson2024}

Conviene mantener una distinción que impida caricaturizar las dos perspectivas.
La regeneración hace posible la vida cristiana y su perseverancia hasta la muerte; el estado intermedio describe la comunión con Cristo de quienes ya han terminado su carrera terrenal; la resurrección del cuerpo consumará la redención de toda la persona.
La lectura aquí defendida no considera falsa la interpretación agustiniana de la vivificación presente, sino insuficiente para identificar el referente inmediato de la visión de Juan.
No todo uso bíblico de \enquote{resurrección} se refiere a la misma experiencia, aunque todas las formas de vida nueva dependan de la victoria histórica del Resucitado.

\section{Las almas que viven y reinan: la paradoja de la victoria}

La primera observación textual es sencilla, pero decisiva: Juan ve \enquote{las almas} de los decapitados por el testimonio de Jesús (Ap 20:4).
La palabra no constituye, por sí sola, una prueba metafísica de incorporeidad, pues el vocabulario bíblico admite usos variados.
Su fuerza surge del contexto visionario.
En Apocalipsis 6:9--11, Juan había visto bajo el altar celestial las almas de quienes fueron muertos por la palabra de Dios; aquellas almas claman, reciben vestiduras blancas y esperan todavía la consumación del juicio.
Son muertos desde el punto de vista terrestre y, sin embargo, están conscientes y son vindicados en la presencia de Dios.
La semejanza entre las dos escenas es demasiado estrecha para ser accidental: testimonio, muerte, almas, vindicación y espera del desenlace final.
Kline subrayó este paralelismo al responder a sus críticos, precisamente porque Apocalipsis 6 presenta sin ambigüedad una situación anterior a la resurrección corporal. \autocite[110--119]{kline1976}

Hill reconoce que los entronizados de 20:4a y los testigos de 20:4b probablemente constituyen un mismo grupo.
Observa, además, que el griego \emph{ezēsan} (ἔζησαν) puede entenderse como \enquote{vivieron} sin que el aoristo imponga la idea ingresiva de \enquote{volvieron corporalmente a la vida}.
De hecho, las almas están vivas cuando se presentan a la mirada del profeta; la visión no narra la apertura de sepulcros ni la restitución de sus cuerpos.
Hill relaciona esta realidad con los vencedores que ya aparecen ante Dios en Apocalipsis 15:2 y con los que \enquote{mueren en el Señor} en 14:13. \autocite[53--54 del PDF del curso]{hill2016}
Esta observación debe formularse con exactitud: la forma verbal admite discusión exegética y no determina ella sola el sentido del pasaje.
Es la convergencia de sujeto, ubicación y paralelos internos la que favorece el referente celestial.

La misma estructura de promesa recorre las cartas a las iglesias.
Al fiel de Esmirna se le manda perseverar hasta la muerte, con la promesa de la corona de vida y de la inmunidad frente a la segunda muerte (Ap 2:10--11).
Los vencedores participan de la autoridad de Cristo (2:26--27; 3:21); aquellos que mueren en el Señor son declarados bienaventurados y descansan de sus trabajos (14:13).
Más adelante, se afirma que los vencedores han triunfado sobre el dragón mediante la sangre del Cordero y el testimonio de su palabra, aun cuando no amaron sus vidas hasta la muerte (12:11).
Apocalipsis 20:4--6 reúne precisamente estos motivos: fidelidad, muerte, vida, reinado, sacerdocio y protección ante la segunda muerte.
La interpretación del estado intermedio nace así del tejido literario del libro y no de una imposición ajena al texto.
Johnson señala, además, que los tronos remiten a la corte celestial de Daniel 7:9--10 y a la escena del trono de Apocalipsis 4; nada obliga a trasladar el reinado a una tierra todavía no renovada. \autocite{johnson2024}

La identidad de los entronizados merece una precisión.
Juan describe de manera sobresaliente a los mártires, porque en ellos la inversión paradójica de la visión alcanza su intensidad máxima.
No obstante, el propio Apocalipsis extiende la bienaventuranza a \enquote{los muertos que mueren en el Señor} (14:13), y las promesas de vencer se dirigen a todas las iglesias.
Por eso, interpretar a los mártires como figura representativa de los fieles difuntos resulta plausible, aunque no debe confundirse con una afirmación explícita de que todos fueron decapitados.
La sangre del testigo no produce su victoria; es la victoria previa del Cordero la que le permite perseverar hasta el final.

Queda por explicar el adjetivo \enquote{primera}.
Kline encontró en Apocalipsis 21:1--5 una clave que desborda la mera secuencia numérica.
Allí el \enquote{primer} cielo y la \enquote{primera} tierra no son simplemente dos objetos anteriores en una lista de otros semejantes; pertenecen al orden provisional que será sustituido por la creación nueva.
La \enquote{segunda muerte}, identificada con el lago de fuego (20:14; 21:8), tampoco es una repetición de la muerte corporal bajo las mismas condiciones.
La relación entre \enquote{primero} y \enquote{segundo} puede expresar, por tanto, una oposición cualitativa entre el orden presente y el definitivo. \autocite[366--370]{kline1975}

Desde este horizonte, la primera resurrección es la participación verdadera, aunque todavía intermedia, de los fieles muertos en la vida del Cristo exaltado.
Es \enquote{primera} no porque sea una de dos restituciones corporales idénticas, sino porque anticipa en la edad presente la resurrección consumada del último día.
Kline llamó la atención sobre la simetría paradójica: la muerte corporal del creyente lo introduce en la vida, mientras la resurrección del incrédulo no lo libra de la segunda muerte.
La diferencia cualitativa no vuelve irreal ninguna de las dos experiencias.
Al contrario, permite reconocer que el lenguaje apocalíptico nombra la verdadera condición de los santos desde la perspectiva del cielo.

\section{La analogía de la fe y el horizonte confesional}

La lectura propuesta debe comprobarse a la luz del testimonio más amplio de la Escritura.
En 1 Corintios 15:20--26, Cristo resucitado es las primicias; después, en su venida, resucitarán quienes le pertenecen; finalmente llegará la consumación, cuando haya sometido a sus enemigos y destruido la muerte.
El acento de Pablo recae tanto en el reinado presente de Cristo como en la resurrección corporal vinculada con su parusía.
Apocalipsis 20 no tiene por qué introducir entre ambos acontecimientos una nueva etapa terrenal prolongada, durante la cual la muerte seguiría activa después de la resurrección gloriosa de los santos.
Es verdad que el texto paulino, considerado aisladamente, no determina todos los intervalos cronológicos posibles; con todo, su movimiento hacia la derrota del último enemigo favorece una consumación unitaria. \autocite{johnson2024}

La misma convergencia aparece en Daniel 12:2, Juan 5:28--29 y Hechos 24:15, pasajes que anuncian la resurrección de justos e injustos bajo el horizonte del juicio definitivo.
Conviene evitar una inferencia excesiva: las formulaciones resumidas no constituyen por sí solas una demostración lógica de simultaneidad instantánea.
Pero sí establecen un patrón canónico que una lectura de dos resurrecciones corporales distantes debe justificar, en lugar de presuponer.
La interpretación intermedia permite leer Apocalipsis 20:4--6 sin deshacer ese patrón: los santos reinan ahora con Cristo y esperan aún la redención de sus cuerpos.

Las confesiones clásicas de la Reforma sostienen precisamente esa doble verdad.
La Confesión de Westminster afirma que las almas de los justos son recibidas en el cielo, donde esperan la redención completa de sus cuerpos; añade que en el último día todos los muertos resucitarán y distingue los destinos de justos e injustos (32.1--3). \autocite{westminster1647}
La Segunda Confesión Bautista de Londres declara igualmente que los creyentes fallecidos son recibidos en el paraíso con Cristo, a la espera de la resurrección corporal de todos en el último día (31.1--3). \autocite{baptist1689}
Ninguno de esos textos define expresamente el significado de la \enquote{primera resurrección}; sería impropio atribuirles la tesis particular de Kline.
Su valor radica en ofrecer una gramática escatológica dentro de la cual la vida consciente con Cristo después de la muerte y la futura resurrección universal no rivalizan entre sí.

La relación entre ambos estados es esencial para evitar un espiritualismo ajeno a la esperanza reformada.
La primera resurrección no reemplaza la restauración del cuerpo ni convierte la gloria definitiva en una supervivencia incorpórea.
El estado intermedio es bendito precisamente porque los santos están con Cristo; es todavía provisional porque la creación entera aguarda su liberación.
Así, la escatología no termina en las almas bajo el altar ni en los tronos del milenio: culmina cuando Dios hace nuevas todas las cosas (Ap 21:1--5).
La esperanza del creyente incluye una comunión real con el Señor inmediatamente después de la muerte y la transformación integral que acompañará su venida.

\section{La objeción premilenial y el peso del contexto}

La objeción más importante proviene de la lectura corporal de Apocalipsis 20.
En los versículos 4--5 aparece dos veces el verbo \emph{ezēsan}: los fieles \enquote{vivieron}, mientras \enquote{los demás muertos no vivieron hasta que se cumplieron mil años}.
El argumento premilenial sostiene que una misma forma verbal, en una oposición tan inmediata, debería conservar un mismo referente: si los demás muertos resucitan corporalmente, también lo hacen los santos.
A esto se añade que \emph{anastasis} (ἀνάστασις) designa ordinariamente la resurrección, y que la lectura más directa de la secuencia parece presentar dos grupos de muertos separados por el milenio.
Ladd representa una formulación clásica de esta orientación, cuya fuerza no debe minimizarse. \autocite{ladd1972}

La objeción obliga a responder en el nivel exacto en que se plantea.
No basta alegar que las visiones son simbólicas.
El propio texto llama \enquote{primera resurrección} a la vida y el reinado de quienes murieron por su testimonio, pero nunca denomina literalmente \enquote{segunda resurrección} a la escena posterior.
Tampoco asegura que \enquote{los demás muertos} reciban vida bienaventurada después de los mil años: aparecen finalmente ante el juicio y la segunda muerte (20:11--15).
La reiteración del verbo admite un contraste deliberado de destinos, no necesariamente dos procesos fisiológicos idénticos.
En la misma unidad literaria, el autor usa \enquote{muerte} para la terminación de la vida corporal y para el castigo eterno; insistir en uniformidad estricta para el verbo \enquote{vivir} sería seleccionar de antemano la literalidad que se pretende demostrar.

Michaels reconoció la relevancia del contraste cualitativo que Kline dedujo de los términos \enquote{primero} y \enquote{segundo}, aunque discutió la reconstrucción de las dos muertes y las dos resurrecciones.
Su respuesta recuerda que la tesis de Kline no debe presentarse como incontrovertible por razones puramente léxicas; la estructura de la visión sigue siendo objeto de disputa. \autocite{michaels1976,kline1976}
Hill, por su parte, ofrece una defensa que no descansa solamente en el adjetivo.
La descripción de almas, la correspondencia con Apocalipsis 6:9--11 y 15:2, y la ausencia de una resurrección corporal narrada hacen que \enquote{vivieron} se entienda mejor como la vida celestial de quienes murieron en el Señor. \autocite[53--54 del PDF del curso]{hill2016}

También hay que considerar el alcance de \enquote{hasta que} (20:5).
La expresión establece el límite temporal de la condición descrita, pero no siempre obliga a que después ocurra la acción contraria con idéntico sentido.
En este pasaje, el destino de los demás muertos se explica a continuación mediante la escena del gran trono blanco; el resultado de su comparecencia no es un segundo reinado, sino el juicio.
En cambio, la suerte de los santos ya ha sido definida por la bienaventuranza del versículo 6 y por la impotencia de la segunda muerte sobre ellos.
La oposición principal de la visión no es entre dos edades de vida corporal terrenal, sino entre la vindicación presente de los fieles y el destino final de quienes permanecieron bajo el señorío de la bestia.

La lectura premilenial conserva una ventaja intuitiva: atribuir a \enquote{resurrección} su sentido corporal más frecuente.
Pero una interpretación responsable no se decide por la frecuencia abstracta de una palabra, sino por su función en una visión determinada.
La lectura del estado intermedio explica conjuntamente por qué Juan ve almas, por qué esas almas pertenecen a testigos ya muertos, por qué viven y reinan antes del juicio definitivo, por qué están protegidas de la segunda muerte y por qué el libro contrapone lo \enquote{primero} a la nueva creación.
Por esa capacidad de integrar los datos sin negar la futura resurrección del cuerpo, constituye una lectura más persuasiva de Apocalipsis 20:4--6.
